un\section{5. Differentialrechnung}

\subsection{Konzept}

\Def $f: D \to \RR$, $x \mapsto y = f(x)$, $D \neq \varnothing$ enthält keine
isolierten Punkte (i.e. Jedes $x \in D$ ist Häufungspunkt von $D \backslash \{x\}$).
\key{mittlere}/\key{durchschnittliche Änderungsrate} oder \key{Differenzenquotient}
von $f$ im Definitionsbereich zwischen $a$ und $a+h$:
$$\frac{\Delta y}{\Delta x} = \frac{f(a+h) - f(a)}{h}$$

\Def Die \key{Ableitung}/\key{Differentialquotient} einer Funkt. $y = f(x)$ \key{an der Stelle $a$}
ist (sofern es existiert):
$$\lim\limits_{h\to 0} \frac{f(a+h)-f(a)}{h} = \lim\limits_{b\to a}\frac{f(b) - f(a)}{b - a} = \left.\frac{dy}{dx}\right|_{x=a} = \left.\frac{df}{dx}\right|_{x=a} = f'(a)$$
Ist \key{momentane Änderungsrate}, Steigung der Tangente,  Zu-/Abnahme

\Def $f: \DD(f) \to \RR$, $x \mapsto y = f(x)$. Die \key{Ableitungsfunktion} ist:
$$a \mapsto f'(a) = \left. \frac{dy}{dx} \right|_{x=a}$$

\Def Funktion $f$. Falls $f'(x)$ existiert, heisst $f$ \key{an der Stelle $x$ differenzierbar}.
Fallst für alle $x \in \DD(f)$ differenzierbar, $f$ \key{differenzierbar}

\Satz $f$ an der Stelle $x_0$ differenzierbar $\implies$ $f$ an der Stelle $x_0$ stetig

\Def $f$ \key{von rechts differenzierbar}, falls ex. \key{rechtsseitige Ableitung}:
$\lim\limits_{\substack{h > x_0 \\ h \to 0}}\frac{f(a+h) - f(a)}{h}$, analog links
$\lim\limits_{\substack{h < x_0 \\ h \to 0}}\frac{f(a+h) - f(a)}{h}$

\Def $f: D \subset \RR \to \RR$. Iterativ \key{höheren Ableitungen} für $n \in \NN_0$:
$f^{(0)} = f, f^{(1)} = f', f^{(2)} = f'', f^{(n+1)} = (f^{(n)})'$.
Falls $f^{(n)}$ existiert, $f$ \key{$n$-fach differenzierbar}. Falls $f^{(n)}$ stetig,
$f$ \key{$n$-fach stetig differenzierbar}.

\Def Menge aller $n$-fach stetig differenzierbarer Funktionen auf D ist $C^n(D)$ und
$C^\infty(D) := \cap_{n=0}^\infty C^n(D)$ sind die \key{glatte Funktionen}

\subsection{Regeln und Sätze}

\begin{minipage}{\linewidth}
\Satz \key{Ableitungsregeln}:

\begin{tabular}{ll}
    $(a \cdot x^p)' = a \cdot p \cdot x^{p-1}$          & $(\cosh(x))' = \sinh(x)$ \\
    $(a^x)' = \ln(a) \cdot a^x \quad \text{für } a > 0$ & $(\ln(x))' = \frac{1}{x}$ \\
    $(\sin(x))' = \cos(x)$                              & $(\arcsin(x))' = \frac{1}{\sqrt{1 - x^2}}$\\
    $(\cos(x))' = -\sin(x)$                             & $(\arccos(x))' = -\frac{1}{\sqrt{1 - x^2}}$ \\
    $(\tan(x))' = \frac{1}{\cos^2(x)} = 1 + \tan^2(x)$  & $(\arctan(x))' = \frac{1}{1 + x^2}$ \\
    $(\sinh(x))' = \cosh(x)$
\end{tabular}
\end{minipage}

\Satz \key{Rechenregeln}: für $f,g: D \to \RR$ an Stelle $x_0$ $n$-fach diff.bar sind
$f+g$ und $f\cdot g$ auch $n$-fach diff.bar. Es gilt
\vspace{-0.6em}
\begin{itemize}
    \item $(f+g)^{(n)}(x_0) = f^{(n)}(x_0) + g^{(n)}(x_0)$
    \item $(f\cdot g)^{(n)}(x_0) = \sum_{k=0}^{n} \begin{pmatrix}n\\k\end{pmatrix} \cdot f^{(k)}(x_0) \cdot g^{(n-k)}(x_0)$
\end{itemize}

\Satz \key{Kettenregel}: $f: D \to E$ an der Stelle $x_0$ diff.bar, $g: E \to \RR$
an Stelle $y_0 = f(x_0)$ diff.bar. $g \circ f: D \to \RR$ an Stelle $x_0$ diff.bar und
$(g \circ f)'(x_0) = g'(f(x_0)) \cdot f'(x_0)$

\Satz \key{Quotientenregel}: $f,g: D \to \RR$ diff.bar an Stelle $x_0$ und $g(x_0) \ne 0$.
$\frac{f}{g}$ in $x_0$ diff.bar und $\left(\frac{f}{g}\right)'(x_0) = \frac{f'(x_0)g(x_0) - f(x_0)g'(x_0)}{g(x_0)^2}$

\Satz \key{Ableitung der Umkehrfunktion}: $f: D \to E \subset \RR$ stetig, bijektiv,  und an Stelle $x_0$ diff.bar,
$f^{-1}: E \to D$, $x_0 \in D$ Häufungspunkt von $D \backslash \{x_0\}$.
Dann $f^{-1}$ an Stelle $y_0 = f(x_0)$ diff.bar: $\left(f^{-1}\right)'(y_0) = \frac{1}{f'(x_0)}$

\Satz \key{Lokale Extremalstellen und erste Ableitung}: $f: D \subset \RR \to \RR$ an lokalen Extremalstelle
$x_0 \in D$ diff.bar, $x_0$ Häufungspunkt von $D \cap (x_0, \infty)$ und $D \cap (-\infty, x_0)$.
Dann gilt $f'(x_0) = 0$

\Satz \key{Lokale Extremalstellen auf Intervall}: $I \subset \RR$ Intervall, $f: I \to \RR$,
$x_0$ lokale Extremalstelle von $f$. Dann mindestens eine der Folgenden zutreffend:
\vspace{-0.6em}
\begin{itemize}
    \item $x_0 \in I$ ist Endpunkt des Intervalls $I$
    \item $f$ an der Stelle $x_0$ nicht differenzierbar
    \item $f$ an der Stelle $x_0$ differenzierbar und $f'(x_0) = 0$
\end{itemize}

\Satz \key{Satz von Rolle}: $a < b$, $f: [a, b] \to \RR$ auf $[a, b]$ stetig und auf $(a, b)$ diff.bar.
Falls $f(a) = f(b)$, dann existiert $\xi \in (a, b)$ mit $f'(\xi) = 0$

\Satz \key{Mittelwertsatz}: $a < b$, $f: [a, b] \to \RR$ auf $[a, b]$ stetig und auf $(a, b)$ diff.bar.
Es gibt ein $\xi \in (a, b)$ mit $f'(\xi) = \frac{f(b) - f(a)}{b - a}$

\Satz \key{Regel von Bernoulli-l'Hôpital}: $a < b$, $f, g: (a, b) \to \RR$, $g, g' \ne 0$, $R > 0$, $h, i: (R, \infty) \to \RR$
\begin{itemize}
    \item $\lim\limits_{\substack{x > a \\ x \to a}} f(x) = \lim\limits_{\substack{x > a \\ x \to a}} g(x) = 0$ oder 
    $\lim\limits_{\substack{x > a \\ x \to a}} |f(x)| = \lim\limits_{\substack{x > a \\ x \to a}} |g(x)| = \infty$
    und $L = \lim\limits_{\substack{x > a \\ x \to a}} \frac{f'(x)}{g'(x)}$ existiert,
    dann gilt $\lim\limits_{\substack{x > a \\ x \to a}} \frac{f(x)}{g(x)} = L$
    \item $\lim\limits_{x\to\infty} f(x) = \lim\limits_{x\to\infty} g(x) = 0$ oder
    $\lim\limits_{x\to\infty} |f(x)| = \lim\limits_{x\to\infty} |g(x)| = \infty$
    und $L = \lim\limits_{x\to\infty} \frac{f'(x)}{g'(x)}$ existiert,
    dann gilt $\lim\limits_{x\to\infty} \frac{f(x)}{g(x)} = L$
\end{itemize}

\subsection{Monotonie, Konvexität und Konkavität}

\Satz $I \subset \RR$ Intervall, $f: I \to \RR$ diff.bar.
$f$ wachsend $\iff$ $f' \ge 0$

\Kor $I \subset \RR$ Intervall, $f: I \to \RR$.
$f$ konst. $\iff$ $f$ diff.bar $\land$ $f' = 0$

\Satz $I \subset \RR$ Intervall, $f: I \to \RR$ diff.bar.
$f$ konvex $\iff$ $f'$ wachsend

\Kor $I \subset \RR$ Intervall, $f: I \to \RR$ 2x diff.bar.
$f$ konvex $\iff$ $f'' \ge 0$

\subsection{Taylorreihen \texorpdfstring{$\star$}{*}}

\Def $n \in \NN_0$, $i \subset \RR$ offenes Int., $f \in C^n(I)$. \key{Taylor-Polynom $n$-ter Ordnung} von $f$
an Stelle $x_0$ ist
$$T_n(x, x_0) = \sum_{k=0}^{n}\frac{f^{(k)}(x_0)}{k!}(x-x_0)^k$$ 

\Satz \key{Restglied von Lagrange}: $n \in \NN_0$, $i \subset \RR$ offenes Int., $f \in C^{n+1}(I)$. Für alle $x \in I$
existiert \key{Zwischenwert} $\xi \in [x_0, x]$ so dass für $R_n := f(x) - T_n(x, x_0)$ gilt:
$$R_n(x) = \frac{f^{(n+1)}(\xi)}{(n+1)!}(x-x_0)^{n+1}$$

\Def $i \subset \RR$ offenes Int., $f \in C^\infty(I)$. \key{Taylorreihe} ist 
$T_\infty(x, x_0) = \sum_{k=0}^{\infty}\frac{f^{(k)}(x_0)}{k!}(x-x_0)^k$
und $f$ heisst \key{analytisch auf $I$} ($f \in C^\omega(I)$) falls $T_\infty(x, x_0)$ für jedes
$x \in I$ gegen $f(x)$ konvergiert ($\lim\limits_{n\to\infty} R_n(x) = 0$)

\Def Taylorreihe mit $x_0 = 0$ heisst auch $\key{Maclaurin-Reihe}$

\Satz \key{Wichtige Maclaurin-Reihen}: Es gilt:
\vspace{-0.6em}
\begin{itemize}
    \item $e^x = \sum_{k=0}^\infty \frac{x^{k}}{k!} = 1 + x + \frac{x^2}{2} + \frac{x^3}{6} + \cdots$
    \item $\cos x = \sum_{k=0}^\infty (-1)^k \frac{x^{2k}}{(2k)!} = 1 + \frac{x^2}{2} + \frac{x^4}{24} - \cdots$
    \item $\sin x = \sum_{k=0}^\infty (-1)^k \frac{x^{2k + 1}}{(2k+1)!} = x + \frac{x^3}{6} + \frac{x^5}{120} - \cdots$
    \item $\ln (1 + x) = \sum_{k=1}^\infty (-1)^{k-1} \frac{x^k}{k} = x - \frac{x^2}{2} + \frac{x^3}{3} - \cdots$ für $|x| < 1$
    \item $\arctan(x) = \sum_{k=0}^\infty (-1)^k \frac{x^{2k+1}}{2k+1} = x - \frac{x^3}{3} + \frac{x^5}{5} - \cdots$ für $|x| \le 1$
    \item $\frac{1}{1-x} = \sum_{k=0}^\infty x^k$ für $|x| < 1$
    \item $(1 + x)^\alpha = \sum_{k=0}^\infty \begin{pmatrix}\alpha \\ k\end{pmatrix} x^k$ für $|x| < 1$, wobei $\begin{pmatrix}\alpha \\ k\end{pmatrix} := \frac{\alpha(\alpha-1)\cdots(\alpha-k+1)}{k!}$
\end{itemize}
